\section{Sequences and Series}\label{sec:series}

There are many kinds of series and sequences, but we care about two of them: \emph{arithmetic} and \emph{geometric}.
In general, a sequence is an ordered list of numbers, and a series is a sum of a sequence.
\subsection{Arithmetic Series}

An arithmetic sequence is one in which the distance between consecutive points are the same:

    \begin{center}
    \begin{tabular}{c|ccccc}
    $i$      & 1 & 2 & 3 & 4 & \dots\\
    \hline
    $x_i$ & -6 & 1 & 8 & 16 & \dots
    \end{tabular}
    \end{center}
In the above, each term has a gap of $7$, which can be worked out by calculating $x_n - x_{n-1}$, and verifying that it is $7$ in each case.

What if we wanted to add the terms together? What if we wanted to find a formula for the arithmetic series? 

A general arithmetic series with $n$ terms looks like the following:
\[ a + (a+d) + (a + 2d) + \dots + (a + (n-1)d)\]

Here $a$ is the first term, and $d$ is the distance between consecutive terms.

Since we are mathematicians, let's give the sum a label $S$.

\[ a + (a+d) + (a + 2d) + \dots + (a + (n-1)d)\]

Here comes a fun trick, if we write the sum backwards:


\[ S = (a + (n-1)d) + (a + (n-2)d) +(a + (n-3)d) + \dots + a\]

Why would we do this? It turns out that adding together the two simplifies our lives greatly:

\begin{align*}
    S + S =& &a &+ &(a+d) &+ &(a+2d)a &+ &\dots &+ &(a + (n-1)d)\\
        &+ &(a + (n-1)d)  &+ &(a + (n-2)d)  &+&(a + (n-3)d) &+ &\dots& + &a\\
        =& &(2a + (n-1)d)  &+ &(2a + (n-1)d)   &+&(2a + (n-1)d)  &+ &\dots& + &(2a + (n-1)d) \\
        =& n(2a + (n-1)d) 
\end{align*}

The last line is since we had $n$ of the same term added together. Therefore we have found:

\[2S = n(2a + (n-1)d) \]

And dividing gives us the formula for an arithmetic series of $n$ terms, starting at $a$ and with a distance of $d$:
\[S = \frac{n}{2}(2a + (n-1)d) \]

\begin{example}
    An arithmetic series is given by

    \[ -6 + 1 + 8 + 15 + 22 + \dots\]

    What is the sum of the first $7$ terms?
\end{example}

\begin{solution}
    If this is an arithmetic series, then each consecutive term must have the same distance apart.

    We can calculate that this distance should be $1 - (-6) = 7$, and can verify this is true for the remaining entries.

    Then all we need to do is plug into our formula, using that the first term is $a = -6$ and we found that $d = 7$, and we want the first $n= 7$ terms:

    \begin{align*}
        S &= \frac{n}{2}(2a + (n-1)d) \\
         &= \frac{7}{2}(2\cdot(-6) + (7-1)7)\\
         &=  \frac{7}{2}(-12 + 6\cdot 7)\\
         &= 7(-6+3\cdot7)\\
         &= 7(15) \\
         &= 105
    \end{align*}

    This finishes the question, finding the sum of the first $7$ terms to be $105$.
    \qed
\end{solution}


\subsection{Geometric Series}

A geometric sequence is one in which the ratio between consecutive points are the same:

    \begin{center}
    \begin{tabular}{c|ccccc}
    $i$      & 1 & 2 & 3 & 4 & \dots\\
    \hline
    $x_i$ & 32 & 16 & 8 & 4 & \dots
    \end{tabular}
    \end{center}
In the above, each term is $1/2$ of the previous, which can be worked out by calculating $x_n / x_{n-1}$, and verifying that it is $1/2$ in each case.

What if we wanted to add the terms together? What if we wanted to find a formula for the geometric series? 

A general arithmetic series with $n$ terms looks like the following:
\[ a + ra + r^2a + \dots + r^{n-2}a + r^{n-1}a\]

Here $a$ is the first term, and $r$ is the ratio between consecutive terms.

Since we are mathematicians, let's give the sum a label $S$.

\[  a + ra + r^2a + \dots + r^{n-2}a + r^{n-1}a\]

Here comes another fun trick, if we write the sum multiplied by $r$:


\[ rS =  ra + r^2a + r^3a + \dots + r^{n-1}a + r^na\]

Why would we do this? It turns out that subtracting to two simplifies our lives greatly:

\begin{align*}
    S - rS =& a &+& \cancel{ra} &+& \cancel{r^2a} &+& \dots  &+& \cancel{r^{n-1}a} & &\\
        & &-&\cancel{ra} &-& \cancel{r^2a}  &-& \dots &-& \cancel{r^{n-1}a} &-& r^na\\
        =& a - r^na
\end{align*}

After factorising the left and right side, we find:

\[S(1-r) = a(1 - r^{n})\]

And dividing gives us the formula for geometric series of $n$ terms, starting at $a$ and with a ratio of $r$:
\[S = \frac{a(1 - r^{n})}{1-r}\]

\begin{example}
    A geometric series is given by

    \[ 32 + 16+ 8 + 4 + \dots\]

    What is the sum of the first $5$ terms?
\end{example}

\begin{solution}
    If this is a geometric series, then each consecutive term must have the same ratio.

    We can calculate that this ratio should be $16/32 = 1/2$, and can verify this is true for the remaining entries.

    Then all we need to do is plug into our formula, using that the first term is $a = 32$ and we found that $r = 1/2$, and we want the first $n= 5$ terms:

    \begin{align*}
    S &= \frac{a(1 - r^{n})}{1-r} \\
      &= \frac{32\left(1 - \left(\frac{1}{2}\right)^5\right)}{1-\frac{1}{2}} \\
      &= 2\cdot 32\left(1 - \left(\frac{1}{2}\right)^5\right) \\
      &= 2\cdot 32\left(1 - \frac{1}{32}\right) \\
      &= 2\cdot 32\cdot \frac{31}{32} \\
      &= 2\cdot 31 \\
      &= 62
\end{align*}


    This completes the question, finding the sum of the first $5$ terms to be $62$.
    \qed
\end{solution}





